% \input{\pSections "sec-theory"}

\section{Theory}

\input{\pSections "ssec-latter"}

%%    %%    %%    %%    %%    %%    %%    %%    %%    %%
\subsection{Big Picture}
\begin{frame}\frametitle{Principals}
\begin{table}[htp]
%\caption{default}
\begin{center}
	\begin{tabular}{ccc}
		 %
	\href{https://en.wikipedia.org/wiki/G._I._Taylor}{
	\begin{overpic}[ height = 1.5in ]
		{\pLocalGraphics people/Taylor}
		%\put(85, 5)	{\colorbox{white}{\cite{taylor1950formationI}}}
	\end{overpic}} &
		 %
	\href{https://en.wikipedia.org/wiki/John_von_Neumann}{
	\begin{overpic}[ height = 1.5in ]
		{\pLocalGraphics people/vN}
		%\put(85, 5)	{\colorbox{white}{\cite{taylor1950formationI}}}
	\end{overpic}} &
		 %
	\href{https://en.wikipedia.org/wiki/Leonid_Sedov}{
	\begin{overpic}[ height = 1.5in ]
		{\pLocalGraphics people/Sedov}
		%\put(85, 5)	{\colorbox{white}{\cite{taylor1950formationI}}}
	\end{overpic}} \\
		%
		Taylor & \ \ von Neumann & Sedov \\
		%
		\pause
		\mg{implosion instability} & \ \mg{explosive lensing} & \href{https://content.time.com/time/subscriber/article/0,33009,869245,00.html}{\mg{Sputnik?}}
		%
	\end{tabular}
\end{center}
\label{tab:mug-shots}
\end{table}%
\end{frame}

%
\begin{frame}\frametitle{Computing Energy From Photographs\jumpLittle}
{\LARGE{
\begin{equation}
	\bl{R} = \bl{S(\gamma) \sqrt[5]{\frac{\rd{E}\, t^{2}}{\rho_{0}} }}
\label{eq:one}
\end{equation}
}}
\end{frame}

%%    %%    %%    %%    %%    %%    %%    %%    %%    %%
%\subsection{Literature Survey}%, Seminal Papers}
\begin{frame}\frametitle{Seminal Papers}
%\caption{default}
\begin{table}[htp]
\begin{center}
	\begin{tabular}{l}
		 %
	\href{run:../local/articles/Taylor50a.pdf}{
	\begin{overpic}[ scale = 0.25 ]
		{\pLocalGraphics ss/taylor-01.pdf}
		\href{\urlTaylora}{\put(85, 5)	{\colorbox{gray!20}{\cite{taylor1950formationI}}}}
	\end{overpic}}  \\\hline
		 %
	\ \\
	\href{run:../local/articles/Taylor50b.pdf}{
	\begin{overpic}[ scale = 0.25 ]
		{\pLocalGraphics ss/taylor-02.pdf}
		\href{\urlTaylorb}{\put(85, 5)	{\colorbox{gray!20}{\cite{taylor1950formationII}}}}
	\end{overpic}} 
		 %
	\end{tabular}
\end{center}
\label{tab:taylor-splash}
\end{table}
\end{frame}
%
\begin{frame}\frametitle{Not \yl{Brook Taylor}}
\begin{table}[htp]
%\caption{default}
\begin{center}
\href{https://en.wikipedia.org/wiki/Brook_Taylor}{
	\begin{tabular}{cc}
		\includegraphics[ width = 1.25in ]{ \pLocalGraphics Taylor/brook-01} &
		\includegraphics[ width = 1.19in ]{ \pLocalGraphics Taylor/brook-02} 
	\end{tabular}
}
\end{center}
%\label{tab:label}
\end{table}
\href{https://mathworld.wolfram.com/TaylorSeries.html}{
\begin{equation*}
	f(x) = f(a) + \frac{x-a}{1!}f' + \frac{(x-a)^{2}}{2!}f''+ \cdots +  \frac{(x-a)^{n}}{n!}f^{(n)} + \cdots
%\label{eq:}
\end{equation*}
}
\end{frame}
%
\begin{frame}\frametitle{Fireball Energy}
\center
	\href{\urlTaylora}{
	\begin{overpic}[ scale = 0.5 ]
		{\pLocalGraphics Taylor/Taylor-01.pdf}
		\put(32, -10)	{\colorbox{white}{$\bl{R = S(\gamma) \sqrt[5]{\frac{E\, t^{2}}{\rho_{0}} }}$}}
		\put(65, -20)	{\colorbox{gray!20}{\cite[p. 172]{taylor1950formationI}}}
	\end{overpic}}
\end{frame}
%
\begin{frame}\frametitle{Fireball Photos: Trinity}
\center
		%
	\href{https://homepage.physics.uiowa.edu/~ghowes/teach/phys7729/docs/Taylor50b.pdf\#page=9}{
	\begin{overpic}[ scale = 0.175 ]
		{\pLocalGraphics taylor/p-photos-01}
		%\put(50, 50)	{\colorbox{white}{$a+b$}}
	\end{overpic}} 
		%
	\href{https://homepage.physics.uiowa.edu/~ghowes/teach/phys7729/docs/Taylor50b.pdf\#page=2}{
	\begin{overpic}[ scale = 0.385 ]
		{\pLocalGraphics taylor/p-table-1}
		%\put(50, 50)	{\colorbox{white}{$a+b$}}
	\end{overpic}}
		%
\end{frame}

\begin{frame}\frametitle{\href{https://homepage.physics.uiowa.edu/~ghowes/teach/phys7729/docs/Taylor50b.pdf\#page=3}{Taylor's Original Figure 1}}
\center
	\href{https://homepage.physics.uiowa.edu/~ghowes/teach/phys7729/docs/Taylor50b.pdf\#page=3}{
	\begin{overpic}[ scale = 0.425 ]
		{\pLocalGraphics taylor/fig-01}
		\pause
		\put(40, 40)	{\colorbox{white}{Method of Choice: \bl{Least Squares}}}
	\end{overpic}}
\end{frame}

%
\begin{frame}\frametitle{Validity of Similarity Relationship}
\center
	\href{https://homepage.physics.uiowa.edu/~ghowes/teach/phys7729/docs/Taylor50b.pdf\#page=3}{
	\begin{overpic}[ scale = 0.5 ]
		{\pLocalGraphics taylor/taylor-basic}
		\put(-45, 105)	{\colorbox{white}{$\frac{5}{2}\log R(t) = (\bl{\solnb} \pm \bl{\errb})\log t + (\solna \pm \erra)$}}
		\put(-15, 40)	{\rotatebox{90}{\colorbox{white}{$\frac{5}{2}\log R$}}}
		\put(37, -8)	{\colorbox{white}{$\log t$}}
	\end{overpic}}
\end{frame}
%\endinput  %  ==  ==  ==  ==  ==  ==  ==  ==  ==
%
\begin{frame}\frametitle{Taylor's Results For Trinity Explosion}
\begin{table}[htp]
%\caption{default}
\begin{center}
	\begin{tabular}{cc}
		%
		Yield & Fireball Rise\\[10pt]
		%
		$E=16.8\ Kt$ & $U=35.7\ m/s$ \\[15pt]
		\pause
		%
	\href{https://homepage.physics.uiowa.edu/~ghowes/teach/phys7729/docs/Taylor50b.pdf\#page=6}{
	\begin{overpic}[ scale = 0.3 ]
		{\pLocalGraphics taylor/p-table-3}
		%\put(50, 50)	{\colorbox{white}{$a+b$}}
	\end{overpic}}  
		&
	\raisebox{0.35cm}{\href{https://homepage.physics.uiowa.edu/~ghowes/teach/phys7729/docs/Taylor50b.pdf\#page=15}{
	\begin{overpic}[ scale = 0.35 ]
		{\pLocalGraphics taylor/p-eq-32}
		%\put(50, 50)	{\colorbox{white}{$a+b$}}
	\end{overpic}}}
		%
	\end{tabular}
\end{center}
%\label{tab:label}
\end{table}
\end{frame}
%
\begin{frame}\frametitle{\href{https://archive.org/details/sedov-similarity-and-dimensional-methods-in-mechanics}{Sedov's Book}}
\center
	\href{https://www.taylorfrancis.com/books/mono/10.1201/9780203739730/similarity-dimensional-methods-mechanics-sedov}{
	\begin{overpic}[ scale = 0.2 ]
		{\pLocalGraphics ss/sedov}
		%\put(85, 5)	{\colorbox{white}{\cite{bethe1944shock}}}
	\end{overpic}}
	\ \\[10pt]
	\cite{sedov2018similarity}
\end{frame}

\begin{frame}\frametitle{Los Alamos: \href{https://www.osti.gov/servlets/purl/4434914}{Shock Hydrodynamics and Blast Wave}}
\center
	\href{run:../local/articles/shock-hydrodynamics-and-blast-waves.pdf}{
	\begin{overpic}[ scale = 0.45 ]
		{\pLocalGraphics ss/bethe-01}
		%\put(85, 5)	{\colorbox{white}{\cite{bethe1944shock}}}
	\end{overpic}}
	\ \\[10pt]
	\href{https://www.osti.gov/servlets/purl/4434914}{\cite{bethe1944shock}}
\end{frame}
%
\begin{frame}\frametitle{Los Alamos \href{https://www.osti.gov/servlets/purl/4326276}{Blast Wave}}
\center
	\href{run:../articles/local/blast-wave.pdf}{
	\begin{overpic}[ scale = 0.25 ]
		{\pLocalGraphics ss/bethe-02}
		%\put(85, 15)	{\colorbox{white}{\cite{bethe1947blast}}}
	\end{overpic}}
	\ \\[10pt]
	\href{https://www.osti.gov/biblio/4326276}{\cite{bethe1947blast}}
\end{frame}
%
\begin{frame}\frametitle{\href{https://pubs.aip.org/aip/jap}{Journal of Applied Physics}}
\center
	\href{run:../local/articles/calculation-hydro-shocks.pdf}{
	\begin{overpic}[ scale = 0.3 ]
		{\pLocalGraphics ss/princeton-01}
		%\put(85, 5)	{\colorbox{white}{\cite{bethe1944shock}}}
	\end{overpic}}
	\ \\[10pt]
	\href{https://doi.org/10.1063/1.1699639}{\cite{vonneumann1950method}}
\end{frame}
%\subsection{RAND}
%
\begin{frame}\frametitle{\href{https://pubs.aip.org/aip/jap}{Journal of Applied Physics}}
\center
	\href{run:../local/article/spherical-blast-waves.pdf}{
	\begin{overpic}[ scale = 0.25 ]
		{\pLocalGraphics ss/brode-01}
		%\put(85, 5)	{\colorbox{white}{\cite{bethe1944shock}}}
	\end{overpic}}
	\ \\[10pt]
	\href{https://doi.org/10.1063/1.1722085}{\cite{brode1955numerical}}
\end{frame}
%
\begin{frame}\frametitle{\href{https://pubs.aip.org/aip/pof}{Physics of Fluids}}
\center
	\href{run:../local/articles/blast-spherical.pdf}{
	\begin{overpic}[ scale = 0.25 ]
		{\pLocalGraphics ss/brode-fluids}
		%\put(85, 5)	{\colorbox{white}{\cite{bethe1944shock}}}
	\end{overpic}}
	\ \\[10pt]
	\href{https://doi.org/10.1063/1.1705911}{\cite{brode1959blast}}
\end{frame}
%
%\begin{frame}\frametitle{\href{https://pubs.aip.org/aip/jap}{Journal of Applied Physics}}
%\center
%	\href{https://pubs.aip.org/aip/jap/article-abstract/19/7/670/159040/The-Attenuation-of-Spherical-Shock-Waves-in-Air}{
%	\begin{overpic}[ scale = 0.4 ]
%		{\pLocalGraphics ss/princeton}
%		%\put(85, 5)	{\colorbox{white}{\cite{bethe1944shock}}}
%	\end{overpic}}
%	\ \\[10pt]
%	\cite{stoner1948attenuation}
%\end{frame}
%Journal of Applied Physics
%
\begin{frame}\frametitle{Los Alamos Technical Report}
\center
	\href{https://cococubed.com/papers/kamm_2000.pdf}{
	\begin{overpic}[ scale = 0.45 ]
		{\pLocalGraphics kamm/kamm}
		%\put(85, 15)	{\colorbox{white}{\cite{bethe1947blast}}}
	\end{overpic}}
	\ \\[10pt]
	\cite{kamm2000evaluation}
\end{frame}

\begin{frame}\frametitle{Wiki: \WikiTvNS s}
\center
	%\href{run:../local/articles/blast-spherical.pdf}{
	%\href{run:../local/articles/Taylor–von-Neumann–Sedov-blast-wave-Wikipedia.pdf}{
	\href{run:../local/articles/Taylor.pdf}{
	\begin{overpic}[ scale = 0.25 ]
		{\pLocalGraphics wiki/splash}
		%\put(85, 5)	{\colorbox{white}{\cite{bethe1944shock}}}
	\end{overpic}}
	\ \\[10pt]
	\href{https://upload.wikimedia.org/wikipedia/commons/c/c1/TNS_blast_wave.pdf}{
	\begin{overpic}[ scale = 0.325 ]
		{\pLocalGraphics plots/TNS_blast_wave}
		%\put(85, 5)	{\colorbox{white}{\cite{bethe1944shock}}}
	\end{overpic}}
\end{frame}

%
\begin{frame}\frametitle{\href{https://apps.dtic.mil/sti/pdfs/ADB178624.pdf\#page=2}{Guide to Atmospheric Tests}\jumpLittle}
\center
	\href{run:../local/articles/ADB178624.pdf}{
	\begin{overpic}[ scale = 0.4 ]
		{\pLocalGraphics ss/guide-atmos-test}
		%\put(85, 5)	{\colorbox{white}{\cite{bethe1944shock}}}
	\end{overpic}}
	\ \\[5pt]
	\href{https://apps.dtic.mil/sti/pdfs/ADB178624.pdf}{\colorbox{gray!10}{\cite{jackson1993guide}}}
\end{frame}
%

%\begin{frame}\frametitle{Taylor's Result Updated}
%\center
%	\href{https://homepage.physics.uiowa.edu/~ghowes/teach/phys7729/docs/Taylor50b.pdf\#page=2}{
%	\begin{overpic}[ scale = 0.55 ]
%		{\pLocalGraphics taylor/fit-result}
%		\put(32, -10)	{\colorbox{white}{\bl{$R(t=1\,s) = 583\,m$}}}
%		\pause
%		\put(20, 35)	{\colorbox{white}{$y = m x + b \quad \Rightarrow \quad y-mx = b$}}
%	\end{overpic}}
%\end{frame}

%%	\bl{y(x)} &= (\solnb \pm \errb)\bl{x} + (\solna \pm \erra) \\[10pt]
%	\bl{\tfrac{5}{2} \log R} &= (\solnb \pm \errb) \bl{\log t} + (\solna \pm \erra)
%\begin{frame}\frametitle{Connecting Fit Parameters}
%Faux linearization:
%\end{frame}

%%    %%    %%    %%    %%    %%    %%    %%    %%    %%
\subsection{Similarity}
%
\begin{frame}\frametitle{Pendulum}
	%\input{\pLocalFigures "fig-pendulumB"}
\end{frame}
%
%%    %%    %%    %%    %%    %%    %%    %%    %%    %%
\subsection{Navier-Stokes Equations}

\begin{frame}\frametitle{Dimensionless Solution Parameters}
\href{\urlTaylora\#page=2}{
\bl{pressure}, \\
\begin{equation}
	y = \frac{p}{p_{0}} = R^{-3} f_{1}
%\label{eq:}
\end{equation}
\bl{density}, \\
\begin{equation}
	\bl{\psi} = \frac{\rho}{\rho_{0}}
%\label{eq:}
\end{equation}
\bl{velocity}, \\
\begin{equation}
	u = R^{-3} \phi_{1}
%\label{eq:}
\end{equation}
}
\end{frame}

\begin{frame}\frametitle{Dimensionless Solution Parameters: \yl{Initial Values}}
At $\eta =1$, 
\ \\
\href{\urlTaylora\#page=4}{
\begin{align}
%
	\psi&= \frac{\gpo} {\gmo} \mg{\ =  6}, \\[10pt]
%
	f &= \frac{2 \gamma} {\gpo} \mg{\ = \frac{7}{6}}, \\[10pt]
%
	\phi &= \frac{2} {\gpo} \mg{\  =\frac{5}{6}} 
%
\end{align}}
% eqns 15a, 16a, 17a
\end{frame}

\begin{frame}\frametitle{Dimensionless Solution Parameters: \yl{Derivatives}}
\href{\urlTaylora\#page=6}{
\begin{align}
%
	f' &= \frac{ (\gamma +6) \eta  \phi -4 \gamma  \phi ^2-6 \eta ^2} {2 \eta  \psi  (\eta -\phi )^2-2 f \eta } f\psi , \\[10pt]
%
	\phi' &= \frac{2 \bl{f'}-3 \gamma  \psi  \phi }{2 \gamma   \psi \paren{\eta - \phi) }}, \\[10pt]
%
	\psi' &= \frac{ \left(\frac{2 \phi }{\eta }+ \bl{\phi'} \right)}{\eta -\phi }\psi 
%
\end{align}}
\end{frame}


\begin{frame}\frametitle{Dimensionless Solution Parameters: Plots}
\center
	\href{\urlTaylorb\#page=5}{
	\begin{overpic}[ scale = 0.75 ]
		{\pLocalGraphics taylor/Taylora-tab01}
		\put(48, -5)	{\colorbox{white}{$\eta$}}
		\put(85, 48)	{\colorbox{white}{$f$}}
		\put(68, 48)	{\colorbox{white}{$\rd{\psi}$}}
		\put(50, 34)	{\colorbox{white}{$\bl{\phi}$}}
	\end{overpic}}
\end{frame}

%
\begin{frame}\frametitle{\WikiEFE}
\center
The \bl{Universe} is an \bl{Ideal Gas}
\end{frame}
%
%
\begin{frame}\frametitle{\WikiEFE}
% Yang, 10.1.54
\href{https://mathworld.wolfram.com/EinsteinFieldEquations.html}{
\begin{equation}
G_{\mu\nu} = - \kappa T_{\mu\nu}
\label{eq:einstein-equations}
\end{equation}
}
\end{frame}

%
\begin{frame}\frametitle{\href{https://mathworld.wolfram.com/TotalDerivative.html}{Total Derivative}}
\begin{equation}
	\frac{D}{Dt} = \partial_{t} + \bl{v^{k}\partial_{k}} = \partial_{t} + \bl{\mathbf{v} \cdot \nabla}	
\label{eq:total derivative}
\end{equation}
\ \\[20pt]
a.k.a \href{https://en.wikipedia.org/wiki/Material_derivative}{Material derivative}, \href{https://mathworld.wolfram.com/ConvectiveDerivative.html}{convective derivative}
%
\end{frame}
%
\begin{frame}\frametitle{Go With The Flow}
\center
	\href{https://math.stackexchange.com/questions/3173298/poincaré-bendixson-theorem-and-limit-cyclefind-the-trapping-region}{
	\begin{overpic}[ scale = 0.25 ]
		{\pLocalGraphics flow/"3173298 Poincare-Bendixson"}
		%\put(50, 50)	{\colorbox{white}{$a+b$}}
	\end{overpic}}
\end{frame}
%
%\begin{frame}\frametitle{Also Know As...}
%\begin{itemize}
%	\item \href{https://bookstore.ams.org/view?ProductCode=CLN/9}{\bl{advective} derivative} (\cite{majda2003introduction})
%	\item \href{https://link.springer.com/book/10.1007/b97537}{\bl{convective} derivative} (\cite{Ockendon_Ockendon_2004})
%	\item \href{https://www.google.com/books/edition/Transport_Phenomena/L5FnNlIaGfcC?hl=en}{derivative \bl{following the motion}} (\cite{bird2002transport})
%	\item \href{https://www.google.com/books/edition/Transport_Phenomena/L5FnNlIaGfcC?hl=en}{\bl{hydrodynamic} derivative}(\cite{bird2002transport})
%	\item \bl{Lagrangian} derivative\footnote{\href{https://www.amazon.com/Introduction-Physical-Oceanography-George-Mellor/dp/1563962101}{Introduction to Physical Oceanography}}
%	\item \href{}{\bl{particle} derivative}
%	\item \href{}{\bl{substantial} derivative}
%	\item \href{}{\bl{substantive} derivative}
%	\item \href{}{\bl{Stokes} derivative}
%	\item \href{https://www.cambridge.org/core/books/gauge-field-theories/BFDECE4B2C0FE428C540232E36D39288}{\bl{Gauge covariant} derivative} (\cite{pokorski2000gauge})
%\end{itemize}
%\end{frame}
%
%\begin{frame}\frametitle{\href{https://en.wikipedia.org/wiki/Material_derivative}{Material Derivative} Also Know As...}
%\begin{enumerate}
%	\item \href{https://bookstore.ams.org/view?ProductCode=CLN/9}{\bl{advective} derivative} (\cite{majda2003introduction})
%	\item \bl{convective} derivative\footnote{\href{https://link.springer.com/book/10.1007/b97537}{Waves and compressible flow}, Ockendon, H. and Ockendon, J. R., Springer, 2004.} 
%	\item \href{https://www.google.com/books/edition/Transport_Phenomena/L5FnNlIaGfcC?hl=en}{derivative \bl{following the motion}} (\cite{bird2002transport})
%	\item \href{https://www.google.com/books/edition/Transport_Phenomena/L5FnNlIaGfcC?hl=en}{\bl{hydrodynamic} derivative} (\cite{bird2002transport})
%	\item \bl{Lagrangian} derivative\footnote{\href{https://link.springer.com/book/9781563962103}{Introduction to Physical Oceanography}, Mellor, G. L., Springer, 1996}
%	\item \href{}{\bl{particle} derivative} (\cite{stoker1992water})
%	\item \href{}{\bl{substantial} derivative} (\cite{bird2002transport})
%	\item \href{}{\bl{substantive} derivative} (\cite{granger1995fluid})
%	\item \href{}{\bl{Stokes} derivative} (\cite{granger1995fluid})
%	\item \href{}{\bl{total} derivative} (\cite{bird2002transport}, \cite{landau1959fluid})
%	\item \href{https://www.cambridge.org/core/books/gauge-field-theories/BFDECE4B2C0FE428C540232E36D39288}{\bl{gauge covariant} derivative} (\cite{pokorski2000gauge})
%\end{enumerate}
%\end{frame}
%
\begin{frame}\frametitle{\href{https://en.wikipedia.org/wiki/Material_derivative}{Material Derivative} Also Know As...}
\begin{enumerate}
	\item \href{https://bookstore.ams.org/view?ProductCode=CLN/9}{\bl{advective} derivative} (\cite{majda2003introduction})
	\item \bl{convective} derivative\footnote{\href{https://link.springer.com/book/10.1007/b97537}{Waves and compressible flow}, Ockendon, H. and Ockendon, J. R., Springer, 2004.} 
	\item \href{https://www.google.com/books/edition/Transport_Phenomena/L5FnNlIaGfcC?hl=en}{derivative \bl{following the motion}} (\cite{bird2002transport})
	\item \href{https://www.google.com/books/edition/Transport_Phenomena/L5FnNlIaGfcC?hl=en}{\bl{hydrodynamic} derivative} (\cite{bird2002transport})
	\item \bl{Lagrangian} derivative\footnote{\href{https://link.springer.com/book/9781563962103}{Introduction to Physical Oceanography}, Mellor, G. L., Springer, 1996}
	\seti
\end{enumerate}
\end{frame}
%
\begin{frame}\frametitle{\href{https://en.wikipedia.org/wiki/Material_derivative}{Material Derivative} Also Know As...}
\begin{enumerate}
	\conti
	\item \href{}{\bl{particle} derivative} (\cite{stoker1992water})
	\item \href{}{\bl{substantial} derivative} (\cite{bird2002transport})
	\item \href{}{\bl{substantive} derivative} (\cite{granger1995fluid})
	\item \href{}{\bl{Stokes} derivative} (\cite{granger1995fluid})
	\item \href{}{\bl{total} derivative} (\cite{bird2002transport}, \cite{landau1959fluid})
	\item \href{https://www.cambridge.org/core/books/gauge-field-theories/BFDECE4B2C0FE428C540232E36D39288}{\bl{gauge covariant} derivative} (\cite{pokorski2000gauge})
\end{enumerate}
\end{frame}
%
%
%%    %%    %%    %%    %%    %%    %%    %%    %%    %%
\subsection{Similarity and dimensional methods}

\begin{frame}\frametitle{\href{https://pubs.aip.org/physicstoday}{Physics Today}}
\center
	%\href{https://pubs.aip.org/physicstoday/article-abstract/64/9/42/413713/Dynamic-similarity-the-dimensionless}{
	\href{run:../local/articles/dynamic-similarity.pdf}{
	\begin{overpic}[ scale = 0.45 ]
		{\pLocalGraphics pt-dynamic-similarity}
		%\put(50, 50)	{\colorbox{white}{$a+b$}}
	\end{overpic}}
	\ \\[10pt]
	\cite{bolster2011dynamic}
\end{frame}

%
\begin{frame}\frametitle{\href{https://scienceworld.wolfram.com/physics/BuckinghamsPiTheorem.html}{Buckingham's Pi Theory}}
	%
\center
Physical laws are \bl{independent} of the form of the \bl{units}
	%
\end{frame}
%
%
\begin{frame}\frametitle{Similarity}
\begin{enumerate}
	\item Nuclear Explosions (\cite{taylor1950formationI, taylor1950formationI})
	\item Astrophysics (\cite{ostriker1988astrophysical})
	\item Laser driven (\cite{edwards2001investigation})
	\item Plasma (\cite{moore2005tailored})
	\item Granular (\cite{boudet2013unstable, boudet2009blast})
\end{enumerate}
\end{frame}
%

\begin{frame}\frametitle{\href{https://www.genealogy.math.ndsu.nodak.edu/id.php?id=102065}{Sedov}'s Formal Development }
\center
	\href{https://www.taylorfrancis.com/books/mono/10.1201/9780203739730/similarity-dimensional-methods-mechanics-sedov}{
	\begin{overpic}[ scale = 0.35 ]
		{\pLocalGraphics sedov-toc}
		%\put(50, 50)	{\colorbox{white}{$a+b$}}
	\end{overpic}}
\end{frame}

%%    %%    %%    %%    %%    %%    %%    %%    %%    %%
%\subsection{Taylor's Papers}
%
%%    %%    %%    %%    %%    %%    %%    %%    %%    %%
%\subsection{Similarity and dimensional methods in mechanics}
\begin{frame}\frametitle{Equations of Motion}
\center
	\href{https://archive.org/details/sedov-similarity-and-dimensional-methods-in-mechanics}{
	\begin{overpic}[ scale = 0.35 ]
		{\pLocalGraphics sedov-148}
		%\put(50, 50)	{\colorbox{white}{$a+b$}}
	\end{overpic}}
	\ \\
	\cite{sedov2018similarity}
\end{frame}
%
\begin{frame}\frametitle{Units: Familiar Quantities}
\begin{enumerate}
	\item velocity
	\item momentum
	\item kinetic energy
	\item potential energy (gravitational)
	\item force
	\item pressure
	\item density (volumetric)
\end{enumerate}
\end{frame}
%
\begin{frame}\frametitle{Units: Familiar Quantities}
\begin{table}[htp]
%\caption{default}
\begin{center}
	\begin{tabular}{lll}
			%
		velocity & $v$ & $L T^{-1}$ \\[5pt]
			%
		momentum & $p = mv$ & $M L T^{-1}$ \\[5pt]
			%
		kinetic energy & $K=\tfrac{1}{2} m v^{2}$ & $M L^{2} T^{-2}$ \\[5pt]
			%
		potential energy & $U = mgh$ & $M L^{2} T^{-2}$ \\[5pt]
			%
		force & $F = ma$ & $M L T^{-2}$ \\[5pt]
			%
		pressure & $P = F / A$ & $M T^{-2}$ \\[5pt]
			%
		density & $\delta$ & $M L^{3}$ \\[5pt]
			%
	\end{tabular}
\end{center}
\label{tab:units}
\end{table}%
\end{frame}

%
\begin{frame}\frametitle{Input Parameters with Dimension}
\begin{enumerate}
	\item $E$, energy
	\item $\rho_{0}$, air pressure
\end{enumerate}
\end{frame}
%
%
\begin{frame}\frametitle{\href{https://en.wikipedia.org/wiki/Rankine\%E2\%80\%93Hugoniot_conditions}{Rankine-Hugoniot} Jump Conditions}
\center
	Connects physical states across shock front
\end{frame}
%
\begin{frame}\frametitle{\href{https://en.wikipedia.org/wiki/Rankine\%E2\%80\%93Hugoniot_conditions}{Rankine-Hugoniot} Jump Conditions}
% https://physics.stackexchange.com/questions/611651/why-does-entropy-jump-across-a-shockwave
\center
	Conservation Laws:\\[20pt]
\begin{table}[htp]
%\caption{default}
\begin{center}
\begin{tabular}{lrcl}
	\bl{mass} & $\rho_{1} u_{1}$ & $=$ & $\rho_{2} u_{2}$\\[10pt]
	\bl{momentum} & $\rho_{1} u_{1}^{2} + p_{1}$ & $=$ & $\rho_{2} u_{2}^{2} + p_{2}$\\[10pt]
	\bl{energy} & $h_{1} + \tfrac{1}{2} u_{1}^{2}$ & $=$ & $h_{2} + \tfrac{1}{2} u_{2}^{2}$
\end{tabular}
\end{center}
\label{tab:conserved}
\end{table}%
\ \\
\pause
Not conserved: \rd{\href{https://physics.stackexchange.com/questions/611651/why-does-entropy-jump-across-a-shockwave}{Entropy}}
\end{frame}
%
%
\begin{frame}\frametitle{\href{https://en.wikipedia.org/wiki/Rankine\%E2\%80\%93Hugoniot_conditions}{Rankine-Hugoniot} Jump Conditions}
% https://physics.stackexchange.com/questions/611651/why-does-entropy-jump-across-a-shockwave
\center
	Conservation Laws:\\[20pt]
\begin{table}[htp]
%\caption{default}
\begin{center}
\begin{tabular}{lrclcrcl}
	mass & \mg{$\rho_{1} u_{1}$} & \mg{$=$} & \mg{$\rho_{2} u_{2}$} & $\Rightarrow$ & $u_{k}$ & $=$ & $\frac{m} {\rho_{k}}$\\[10pt]
	mom. & \mg{$\rho_{1} u_{1}^{2} + p_{1}$} & \mg{$=$} & \mg{$\rho_{2} u_{2}^{2} + p_{2}$} & $\Rightarrow$ & $p_{2} - p_{1}$ & $=$ & $m \paren{u_{1} - u_{2}}$\\[10pt]
	energy & \mg{$h_{1} + \tfrac{1}{2} u_{1}^{2}$} & \mg{$=$} & \mg{$h_{2} + \tfrac{1}{2} u_{2}^{2}$} & $\Rightarrow$  & $h_{2} - h_{1}$ & $=$ & $\tfrac{1}{2} \paren{u_{1}^{2} - u_{2}^{2}}$
\end{tabular}
\end{center}
\label{tab:conserved}
\end{table}%
\end{frame}
%
\begin{frame}\frametitle{\href{https://en.wikipedia.org/wiki/Rankine\%E2\%80\%93Hugoniot_conditions}{Rankine-Hugoniot} Jump Conditions}
\center
	Implications:
\begin{equation}
	-m^{2} = \frac{p_{2} - p_{1}} { 1 / \rho_{2} - 1 / \rho_{1} }
\label{eq:Michelson–Rayleigh}
\end{equation}
%
\begin{equation}
	h_{2} - h_{1} = \frac{1}{2} \displaystyle\paren{ 1 / \rho_{2} - 1 / \rho_{1} } \paren{p_{2} - p_{1}}
\label{eq:enthalpy}
\end{equation}
\end{frame}

%\section{Jump Conditions}
\begin{frame}\frametitle{Jump Conditions: Physical Properties}
\begin{table}[htp]
\caption{default}
\begin{center}
\begin{tabular}{lcll}
	$u_{l}$ && $u_{r}$ & density\\
	$\rho_{l}$ && $\rho_{r}$ & velocity\\
	$T_{l}$ && $T_{r}$ & pressure\\
\end{tabular}
\end{center}
\label{tab:}
\end{table}%
\end{frame}

%\section{Shock Hydrodynamics}
%\begin{frame}\frametitle{Shock Hydrodynamics}
%\center
%	\href{http://homepage.physics.uiowa.edu/~ghowes/teach/phys7729/docs/Taylor50a.pdf}{
%	\begin{overpic}[ scale = 0.5 ]
%		{\pLocalGraphics ss/taylor-01}
%		%\put(50, 50)	{\colorbox{white}{$a+b$}}
%	\end{overpic}}
%	\cite[LANL]{bethe1944shock}
%\end{frame}

\begin{frame}\frametitle{Physical Properties}
\begin{table}[htp]
\caption{default}
\begin{center}
\begin{tabular}{lcll}
	$u_{l}$ && $u_{r}$ & density\\
	$\rho_{l}$ && $\rho_{r}$ & velocity\\
	$T_{l}$ && $T_{r}$ & pressure\\
\end{tabular}
\end{center} 
\label{default}
\end{table}%
\end{frame}

\begin{frame}\frametitle{Assumptions, Simplifications}
\begin{enumerate}
	\item point explosion
	\item infinitesimal time
	\item high energy explosion
	\item infinitely strong shock
	\item 
\end{enumerate}
\end{frame}

\begin{frame}\frametitle{Assumptions, Simplifications}
\begin{enumerate}
	\item spherically symmetric blast wave
	\item point source
	\item infinitesimal time
	\item infinitely strong shock
	\item 
\end{enumerate}
\end{frame}

%
\begin{frame}\frametitle{Self-Similar Region}
\begin{enumerate}
	\item pressure behind shock wave dominates air pressure $\rho_{1} >> \rho_0$
	\item point source
	\item infinitesimal time
	\item infinitely strong shock
	\item 
\end{enumerate}
\end{frame}
%


\endinput  %  ==  ==  ==  ==  ==  ==  ==  ==  ==
``''
Buckingham's Pi theorem, also known as the Buckingham π theorem, is a fundamental principle in dimensional analysis and fluid mechanics. It provides a method for determining the relationship between physical quantities in terms of their dimensions. The theorem is particularly useful when dealing with complex physical systems where multiple variables are involved.

The theorem states that if there is a physical relationship between 'n' variables involving 'm' fundamental dimensions (such as length, mass, time, etc.), then the relationship can be expressed in terms of 'n - m' dimensionless groups. These dimensionless groups are formed by combining the original variables in a way that eliminates the dimensions.

Now, let's discuss the connection between Buckingham's Pi theorem and self-similarity.

Self-similarity is a property of certain objects or systems where the structure at a small scale is similar to the structure at a larger scale. In the context of fluid dynamics, self-similarity often arises in the study of turbulent flows and fractals. The idea is that certain statistical properties of the flow remain consistent across different scales.

Buckingham's Pi theorem is related to self-similarity in the sense that it helps identify and express dimensionless groups that capture the essential characteristics of a system, which may include self-similar behavior. When analyzing a self-similar system, dimensional analysis can help reveal dimensionless groups that remain constant across different scales.

For example, in the study of turbulent flows, the Reynolds number (Re) is a dimensionless quantity that characterizes the flow. The Reynolds number is formed using variables such as fluid velocity, characteristic length, and viscosity. In a self-similar turbulent flow, the Reynolds number may remain constant even as the scale of observation changes.

In summary, Buckingham's Pi theorem is a \href{https://mathworld.wolfram.com/Power.html}{Power} ful tool for identifying dimensionless groups in physical relationships, and these dimensionless groups can be crucial in understanding and expressing self-similar behavior in certain systems, such as turbulent flows.
